% Generated draft: requires mathematical and source review.
\documentclass[11pt]{article}
\usepackage[margin=25mm]{geometry}
\usepackage{fontspec}
\setmainfont{Latin Modern Roman}
\usepackage{amsmath,amssymb,mathtools,unicode-math}
\setmathfont{Latin Modern Math}
\usepackage{enumitem,needspace,xurl,hyperref}
\setlength{\parindent}{0pt}
\setlength{\parskip}{5pt}
\setlength{\emergencystretch}{3em}
\newcommand{\N}{\mathbb{N}}
\newcommand{\Z}{\mathbb{Z}}
\newcommand{\Q}{\mathbb{Q}}
\newcommand{\R}{\mathbb{R}}
\newcommand{\C}{\mathbb{C}}
\newcommand{\F}{\mathbb{F}}
\newcommand{\A}{\mathbb{A}}
\newcommand{\PP}{\mathbb{P}}
\newcommand{\E}{\mathbb{E}}
\newcommand{\eps}{\varepsilon}
\newcommand{\dd}{\,\mathrm{d}}
\begin{document}
\section*{UnsolvedMath ID: 6200043}
\textbf{Source problem code:} catalogue entry
\subsection{Definitions and mathematical statement}
Let $G$ be the intended word-hyperbolic group and let $X=\partial_\infty G$ be its boundary equipped with the metric and measure relevant to the Loewner hypothesis. For clarity, a usual $Q$-Loewner framework uses $Q>1$, an Ahlfors $Q$-regular measure $\mu$, and a positive function $\phi$ such that
\[
 \operatorname{Mod}_Q\Gamma(E,F)\ge
 \phi\!\left(\frac{\operatorname{dist}(E,F)}
                    {\min\{\operatorname{diam}E,\operatorname{diam}F\}}\right)
\]
for disjoint nondegenerate continua $E,F\subseteq X$. Here $\Gamma(E,F)$ is the family of joining curves, and
\[
 \operatorname{Mod}_Q\Gamma=\inf_\rho\int_X\rho^Q\,\mathrm d\mu,
 \qquad\int_\gamma\rho\,\mathrm ds\ge1\quad(\gamma\in\Gamma),
\]
with the infimum over admissible nonnegative densities. Ahlfors regularity means $\mu(B(x,r))$ is comparable to $r^Q$ uniformly over the relevant radii. The source does not supply $Q$, the metric, the measure, or an exact Loewner convention.

An embedding $f:X\to X$ is quasisymmetric if it is a homeomorphism onto its image and there is a homeomorphism $\eta:[0,\infty)\to[0,\infty)$ such that for distinct $a,b,x$,
\[
 \frac{d(f(x),f(a))}{d(f(x),f(b))}
       \le\eta\!\left(\frac{d(x,a)}{d(x,b)}\right).
\]
Quasisymmetric cohopficity means that every such self-embedding is surjective.

Heinonen's conjecture is that a Loewner boundary $\partial_\infty G$ is quasisymmetrically cohopfian:
\[
 X\text{ Loewner}\quad\Longrightarrow\quad
 \forall f:X\to X\text{ quasisymmetric embedding},\ f(X)=X.
\]
The source explicitly proposes boundaries of Fuchsian buildings as a test case. It is the no-proper-self-embedding conclusion, not the tautology that a self-homeomorphism is onto, that is asked.
\subsection{Short English statement}
Must every quasisymmetric embedding of a Loewner hyperbolic-group boundary into itself cover the whole boundary? Test the conjecture on Fuchsian-building boundaries.

The supplied research summary says rigidity is known for some Loewner boundaries and gives an informal round-carpet/rank-one-building example, but provides no exact space, theorem, metric, or reference. It calls the general assertion unresolved. The OPEN-TRIAGE label and 2005 source are retained with current status explicitly needing review; the example is not adopted as a verified special-case theorem.
\subsection{Sources}
\begin{itemize}
\item Local UnsolvedMath registry, permanent ID 6200043, source code catalogue entry, ``Boundaries of Groups and Kleinian Groups --- Problem 43,'' proposed by Juha Heinonen, 2005. Dataset curation and metadata: Ulam AI and UnsolvedMath Contributors, CC BY 4.0.
\item Kapovich, \emph{Problems on Boundaries of Groups and Kleinian Groups} (2005), Problem 43, PDF page 13, \url{https://www.math.ucdavis.edu/~kapovich/EPR/problems.pdf}; normalized pdf-text extraction, accessed 2026-07-29.
\item The supplied summary's claimed building/round-carpet special cases have no separate primary reference or precise hypotheses. Source-specific provenance marks current status and statement-reuse terms NEEDS\_REVIEW for the public author-hosted PDF, separately from the dataset licence.
\end{itemize}
\end{document}
